\section{Subsequent Closed-Loop Research Iterations}
The previous pilot is preserved as a versioned record, not replaced by this section. New runs, checkpoints, actual camera recordings and complete case tables are indexed in \texttt{08\_ITERATIVE\_RESEARCH}. Physical task versions and development cohorts are kept separate; results below are not new final-test scores or official benchmark replications.

\begin{table*}[t]\centering\small
\caption{Completed new-round P2 development evaluations generated from the actual run index. Independent success is distinct from raw reward. Seeds refer to independently trained policies, not independent replications of the same test scenes. Each round has a separately versioned physical interface and duration contract, so cross-round differences are not a single-component causal ablation. Camera re-executions, training-geometry checks and incomplete audits are excluded.}\label{tab:iterations}
\begin{tabular}{llrrrr}
\toprule
Round & Method & Seed & Obs./act. & Train interactions & Independent success \\
\midrule
R1 & Feedback & -- & 64/3 & 0 & 25/32 \\
R1 & Parameter PPO & 31 & 64/3 & 32768 & 25/32 \\
R1 & Effect-history PPO & 31 & 64/3 & 32768 & 25/32 \\
R2 & Feedback & -- & 66/4 & 0 & 25/32 \\
R2 & Yaw parameter PPO & 31 & 66/4 & 32768 & 24/32 \\
R2 & Yaw + history PPO & 31 & 66/4 & 32768 & 24/32 \\
R3 & Feedback / nominal & -- & 66/4 & 0 & 26/32 \\
R3 & Re-contact feedback & -- & 87/5 & 0 & 26/32 \\
R3 & Event / ordinary PPO & 31 & 87/5 & 32768 & 22/32 \\
R3 & Event / ordinary PPO & 41 & 87/5 & 32768 & 25/32 \\
R3 & Event / ordinary PPO & 51 & 87/5 & 32768 & 20/32 \\
R3 & Event / active PPO & 31 & 87/5 & 32768 & 24/32 \\
R3 & Event / active PPO & 41 & 87/5 & 32768 & 24/32 \\
R3 & Event / active PPO & 51 & 87/5 & 32768 & 18/32 \\
R3 & Fixed / ordinary PPO & 31 & 87/5 & 32768 & 25/32 \\
R3 & Fixed / active PPO & 31 & 87/5 & 32768 & 25/32 \\
\bottomrule
\end{tabular}

\end{table*}

\subsection{R0: visual re-execution of the existing residual gain}
A fresh complete physical re-execution of the historical seed41 ordinary residual policy and the feedback controller obtains 18 versus 8 accepted completions out of 32 matched geometry cases. The task decomposition is P1: 12 versus 3 out of 16, and P2: 6 versus 5 out of 16. Hence most of the aggregate gain is attributable to P1; it cannot be assigned to a new P2 method. Representative P2 case001 is a feedback failure and residual success, whereas case003 fails for both and case009 succeeds for both. These illustrative videos are selected by outcome type from the full case table, not presented as an unbiased random sample or as a novel-mechanism evaluation. Recorded robot and object initial states match; hidden PhysX solver state is not claimed to be restored.

\safewidefigure{R0_case001_filmstrip.pdf}{Fresh R0 closed-loop PhysX evaluation of historical P2 case001: feedback fails and ordinary residual PPO passes independent qualification. The same physical clocks are sampled in both rows; END/HOLD explicitly denotes a completed episode held for comparison. This is an existing-policy example, not the new candidate method.}{fig:r0paired}

\subsection{R1: actionable contact parameter and effect history}
The first new iteration gives both ordinary parameter PPO and its effect-history variant the same three-dimensional legal interface, including a phase1 contact-target offset bounded by 8 mm. Both receive 32,768 online interactions, the same 64-dimensional capacity, shared reward and physical controller; the effect-history channels are zeroed in the ordinary version. Each run completes 48 training episodes. Independent evaluation on 32 development P2 cases gives 25/32 for both, with identical successful- and failed-case sets and zero recorded tool-contract violations; same-version feedback also achieves 25/32. This development cohort differs from v7 and must not be compared by subtracting its count from v7's 30/32.

Across jointly successful cases the history policy is 0.0373 s slower, using episode control steps divided by 60 Hz rather than global simulation time. Case008 shows real action and contact differences but fails under both policies. A wider allowed parameter is not evidence that useful re-contact was learned. Observed phase1 failures motivate a shared rotated-contact footprint repair and legal tool-yaw exploration, whose engineering effects must be available to all baselines and not attributed to reinforcement-learning novelty.

\safewidefigure{R1_case008_filmstrip.pdf}{R1 case008, actual same-case executions after new online training. Both strategies fail despite different object/contact trajectories. The complete 32-case comparison, rather than this example alone, supports the equal 25/32 outcome.}{fig:r1paired}

\subsection{R2: shared contact geometry repair and finite tool yaw}
The second physical iteration shares a rotated footprint and separating-axis contact geometry repair across all controllers and gives both learning policies the same finite tool-yaw interface (66 observation and four action dimensions). The repair is an engineering correction, not a new RL mechanism. Two additional models each receive 32,768 interactions. Independent development acceptance is 24/32 for both ordinary parameter PPO and effect-history PPO, again with identical outcome sets; shared strong feedback achieves 25/32. History is 0.3083 s slower across jointly successful cases. These are separate task-version evaluations, not an ablation that attributes any R1--R2 difference to a single algorithm component.

\safewidefigure{R2_case005_filmstrip.pdf}{Actual R2 P2 case005 closed-loop comparison after rotated geometry and finite tool-yaw repair. Both policies fail. Their initial recorded physical states match, while contact trajectories differ. The shared engineering changes do not demonstrate an RL method contribution.}{fig:r2paired}

\subsection{R3: physical re-contact and execution-active likelihood}
All three event-refresh seed pairs complete training and independent development audits. Ordinary PPO obtains 22/25/20 accepted successes out of 32, versus execution-active PPO's 24/24/18. Mean success is 69.79\% with seed sample SD 7.86 percentage points versus 68.75\% with SD 10.83 points; the paired count differences are +2/-1/-2. These are three independent training seeds on the same 32 development scenes, not 96 independent scenes. Both same-interface feedback with automatic re-contact and the best nominal feedback without it achieve 26/32. Each active-policy successful-case set is a subset of feedback's, with no rescue and 2/2/8 regressions. Neither improved completion nor reduced seed variability is supported.

For seed31, active PPO rescues three ordinary-PPO failures but regresses on one, a genuine +2-case local signal. It does not replicate across seeds. The all-case mean duration falls 0.7474 s for that seed, but the 21 jointly successful cases are 0.3325 s slower. The corresponding common-success differences are also slower by 0.6632 and 0.3139 s for seeds41/51. Against strong feedback, the active policy is slower on shared successes by 0.227/0.560/0.060 s. Thus selected rescues or a favorable all-case mean cannot establish a general time advantage.

The fixed-refresh seed31 ablation gives 25/32 for both ordinary and active likelihood, with identical accepted-case sets. Compared with event refresh's 22/24, this does not support a benefit of the contact-event heuristic. On the 25 common fixed-refresh successes, the active version is 0.208 s faster. This is retained as a real single-seed secondary-cost signal, not a repeatable primary-success advantage or a cross-seed method conclusion. The four cells share the executor, 24 s contract and training budget; the event/fixed comparison is not mixed with R1/R2's shorter task version.

\safewidefigure{R3_case005_filmstrip.pdf}{Actual seed31 P2 case005, an ordinary-policy failure and execution-active success. This represents one of three rescues and one regression in that seed, not a repeatable three-seed advantage. No learned recovery is inferred: the active training run has no re-contact activation. All case outcomes and the stronger feedback comparison remain reported.}{fig:r3paired}

The new formal training comprises twelve models at 32,768 interactions each, or 393,216 interactions, including both fixed-refresh ablations. Two software/optimizer smokes add 5,120, a separate completed live-training capture adds 12,288, and its retained failed recording attempt adds 4,096 interactions: 414,720 charged online interactions in total. There are 15 completed training checkpoints, not 15 formal comparison seeds. The failed attempt reaches its first rollout then fails in camera-wrapper serialization before an optimizer update; its known work is charged, while unknown reset-related physics cost is not invented. Physical probes and camera evaluations are separately indexed, not new independent training scenes.

\safewidefigure{R3_online_training_16env_montage.pdf}{Actual additional online-training capture: sixteen camera views of independent environments in one Isaac application at the same physics tick1536, not a single-camera panorama or sixteen applications. The run performs 12,288 additional interactions, three rollout updates and 288 optimizer minibatches. It is not one of the fixed-budget formal performance models. Full video, a separate true wide-camera view and before/after optimizer-state evidence are included.}{fig:actualtraining}

\safewidefigure{FigR_online_learning.pdf}{Actual completed R1--R3 online episode records. Rewards and training episode outcomes are not independently accepted development performance. Model seeds and physical task versions are distinguished; software probes and showcase continuation do not inflate the formal curves.}{fig:closedlearning}
\safewidefigure{FigR_development_cases.pdf}{All matched development-case outcomes, including shared failures, rescues and regressions. R1/R2 and R3 use separately versioned interfaces and horizons. The complete case matrix, not a selected filmstrip, supports the method decision.}{fig:closedcases}

\paragraph{Exploratory no-guide pushing transfer.}
Removing the guide rails provides a physically different contact condition without additional policy training. On 32 exploratory cases, feedback and the ordinary seed31 policy each pass 32/32, while the execution-active seed31 policy passes 29/32 and is 2.588 s slower on common successes. This is a zero-training transfer probe, not a sealed final test or evidence of shared-policy improvement. It does not turn the historical P1/P2 results into a verified cross-task benefit of the new mechanism.

\paragraph{Real-rollout actor-gradient diagnostic.}
The formal seed31 ordinary and execution-active runs save exactly equal first 512 observations, sampled actions, GAE values and execution masks, with exactly equal pre-update policy state dictionaries. On these real on-policy samples, a same-weight diagnostic changes the score-gradient vector by Euclidean norm 0.405706 with cosine 0.933645; old log-probability reconstruction error is below $1.2\times10^{-7}$. It consumes no extra environment interactions or optimizer steps. This isolates the score construction on the inspected samples, not the entire first rollout or a performance gain. The diagnostic uses its own common advantage normalization, and is not substituted for actual training minibatch updates.

For that seed's full 32,768 training interactions, the ordinary policy enters automatic re-contact for 576 control samples and has only three valid normal-angle samples; the execution-active policy has neither. Thus the shared executor is implemented, but it is false to claim the active policy learned recovery or meaningful use of every new action dimension. The completed multi-seed physical results, not a nonzero actor gradient, will determine the mechanism decision.
